% Gemaakt door tools/bouw_bundel.py: niet met de hand aanpassen, maar het sjabloon in tools/bundel/.
\documentclass[11pt,twoside]{book}
\usepackage{stijl}
\input{paginas}
\newcommand\VakNaam{<<VAK>> <<NIVEAU>>}
\newcommand\ExamenDuur{<<DUUR>>}
\newcommand\ExamenHulp{<<HULP>>}
\renewcommand\ExamenTab{<<EXAMENTAB>>}
\newif\ifformulekaart\formulekaart<<FORMULEKAART>>
\ifformulekaart\def\FormuleRij{Formules & blz.~\pageref{formulekaart} \\[3pt]}\else\def\FormuleRij{}\fi
\hypersetup{pdftitle={Eindexamen oefenen: <<VAK>> <<NIVEAU>>},pdfauthor={<<SAMENSTELLER>>},
  pdfsubject={Centrale examens <<VAK>> <<NIVEAU>> <<JAREN>> per onderwerp, met de officiele uitwerkingen}}

\begin{document}
\pagestyle{fancy}

% =================== omslag ===================
\thispagestyle{leeg}
\pdfbookmark[-1]{Omslag}{omslag}
\begin{tikzpicture}[remember picture,overlay]
  \definecolor{omslag}{HTML}{12313B}
  \fill[omslag] (current page.south west) rectangle (current page.north east);
  % ruitjes
  \begin{scope}[shift={(current page.south west)}]
    \begin{scope}[white,opacity=.07,line width=.4pt]
      \foreach \i in {0,...,42} {\draw (0,\i*5mm) -- (148mm,\i*5mm);}
      \foreach \i in {0,...,29} {\draw (\i*5mm,0) -- (\i*5mm,210mm);}
    \end{scope}
    % motief per vak (tools/bundel/omslag/<familie>.tex): lijntekeningen, geen foto's
    \begin{scope}[shift={(20mm,128mm)},x=1mm,y=1mm]
      \input{omslag}
    \end{scope}
  \end{scope}
  % tabs aan de rand in de kleuren van de hoofdstukken
  \foreach \k/\c in {<<OMSLAGTABS>>} {
    \fill[\c] ($(current page.north east)+(-4.2mm,-24mm-\k*25mm+25mm)$) rectangle ++(6mm,-21mm);}
  \node[anchor=north west,text=white,font=\sffamily\halfvet\footnotesize] at ($(current page.north west)+(14mm,-14mm)$)
    {\addfontfeature{LetterSpace=12}EINDEXAMEN OEFENEN};
  \node[anchor=north west,text=white,align=left] at ($(current page.north west)+(13mm,-96mm)$)
    {{\sffamily\zwart\fontsize{44}{44}\selectfont <<VAK>>}\\[3mm]
     {\sffamily\licht\fontsize{44}{44}\selectfont\color{kleur2!45!white}<<NIVEAU>>}};
  \node[anchor=north west,text=white,align=left,text width=112mm] at ($(current page.north west)+(14mm,-140mm)$)
    {\sffamily\fontsize{11.5}{15}\selectfont\hyphenpenalty=10000\raggedright Examenopgaven per onderwerp, complete examens en de offici\"ele uitwerkingen};
  \node[anchor=south west,text=white,align=left] at ($(current page.south west)+(14mm,24mm)$)
    {\sffamily\begin{tabular}{@{}l@{\hspace{9mm}}l@{\hspace{9mm}}l@{}}
      {\zwart\fontsize{17}{18}\selectfont <<NOPG>>} & {\zwart\fontsize{17}{18}\selectfont <<NEX>>} & {\zwart\fontsize{17}{18}\selectfont <<JAREN>>}\\[1mm]
      {\footnotesize opgaven per onderwerp} & {\footnotesize complete examens} & {\footnotesize examenjaren}
    \end{tabular}};
  \fill[white,opacity=.12] (current page.south west) rectangle ($(current page.south east)+(0,12mm)$);
  \node[anchor=west,text=white,font=\sffamily\scriptsize] at ($(current page.south west)+(14mm,6mm)$) {Gratis uitgave \textbullet\ niet voor verkoop};
  \node[anchor=east,text=white,font=\sffamily\scriptsize] at ($(current page.south east)+(-14mm,6mm)$) {<<WEBADRES>>};
\end{tikzpicture}
\clearpage
% =================== voorwoord (binnenkant omslag) ===================
\thispagestyle{leeg}
\pdfbookmark[-1]{Voorwoord}{voorwoord}
\vspace*{12mm}
{\sffamily\zwart\fontsize{17}{20}\selectfont Voorwoord}\par\vskip 9mm
{\opmaakproza\fontsize{10.5}{15.5}\selectfont\setlength\parskip{2.6mm}
Beste examenkandidaat,

Voor je ligt een bundel met de centrale examens <<VAKKLEIN>> <<NIVEAU>> van <<JAREN>>. De opgaven zijn
letterlijk overgenomen uit de examens van het College voor Toetsen en Examens, en bij elke vraag hoort het
offici\"ele correctievoorschrift. Je oefent dus met precies het materiaal waarmee je straks wordt beoordeeld.

In deel~1 staan de opgaven uit eerdere examens, geordend naar de domeinen en onderwerpen van de syllabus.
Zo kun je gericht werken aan de onderwerpen die je nog lastig vindt. Deel~2 bevat de <<NEX>> meest recente examens
in hun geheel. Die maak je bij voorkeur aan het einde van je voorbereiding, in de tijd die je ook op het examen hebt.
Achterin vind je de uitwerkingen.

We raden je aan om regelmatig te oefenen, je werk eerlijk na te kijken met het correctievoorschrift en bij te houden
welke onderwerpen nog aandacht vragen. Op de volgende pagina's lees je hoe je dat doet.

Deze bundel is gratis en mag vrij worden gedeeld. We wensen je veel succes met je voorbereiding en met je examen.\par}
\vskip 7mm
{\sffamily\small <<SAMENSTELLER>>\\ \color{grijs}<<MAANDJAAR>>\par}
\vfill
{\opmaakproza\scriptsize\color{grijs}\textbf{Opmerking.} Hoewel deze bundel met de grootst mogelijke zorg is samengesteld,
kunnen de makers geen aansprakelijkheid aanvaarden voor eventuele fouten. Raadpleeg voor actuele informatie over je
examen, zoals de datum, de duur en de toegestane hulpmiddelen, altijd je docent of examenblad.nl.\par}
\clearpage

% =================== titelpagina ===================
\thispagestyle{leeg}
\vspace*{30mm}
{\sffamily\halfvet\footnotesize\color{kleur1}\addfontfeature{LetterSpace=12}EINDEXAMEN OEFENEN}\par\vskip 6mm
{\sffamily\zwart\fontsize{30}{32}\selectfont <<VAK>>\\[1mm]\licht <<NIVEAU>>}\par\vskip 8mm
{\large Opgaven uit de centrale examens <<JAREN>>, per onderwerp\\ met <<NEX>> complete examens en de offici\"ele uitwerkingen\par}
\vfill
{\sffamily\small Een gratis uitgave van <<SAMENSTELLER>>\\ \color{grijs}<<WEBADRES>>\par}
\vskip 4mm
{\sffamily\scriptsize\color{grijs}Met examens en correctievoorschriften van het College voor Toetsen en Examens\par}
\clearpage

% =================== colofon ===================
\thispagestyle{leeg}
\vspace*{\fill}
{\footnotesize\opmaakproza
{\sffamily\bfseries Eindexamen oefenen: <<VAK>> <<NIVEAU>>}\\
Uitgave van <<DATUM>>.

\textbf{Bron.} De examenopgaven, figuren en uitwerkingen in dit boek komen uit de centrale examens en
correctievoorschriften van het College voor Toetsen en Examens (CvTE), gepubliceerd op examenblad.nl.
\textcopyright\ Staat der Nederlanden. Hergebruik is toegestaan met bronvermelding
(examenblad.nl, onder \emph{Over Examenblad}, \emph{Copyright}). De oorspronkelijke bron van elke opgave staat erbij;
de bronnen van foto's en figuren staan achterin (blz.~\pageref{bronnen}). Aanvullingen en errata van het CvTE zijn verwerkt.

\textbf{Niet van het CvTE.} Dit boek is niet gemaakt door of in opdracht van het CvTE. De indeling naar onderwerp,
de uitleg en de opmaak zijn van Eindexamen Oefenen. Heb je een fout gevonden? Meld het via <<WEBADRES>>; daar staat ook wat er na elke uitgave is verbeterd.

\textbf{Gratis.} Dit boek is gratis en mag vrij worden gedeeld en geprint, maar niet worden verkocht.

\textbf{Lettertypes.} Source Serif 4 en Source Sans 3 (Adobe, SIL Open Font License). Gezet met Lua\LaTeX.\par}
\clearpage

% =================== inhoud ===================
\renewcommand\kopL{Inhoud}\renewcommand\kopR{Inhoud}
\phantomsection\label{inhoud}\pdfbookmark[-1]{Inhoud}{inhoud}
{\sffamily\zwart\fontsize{19}{22}\selectfont Inhoud}\par\vskip 8pt
\Inhoud{inhoud}
\clearpage

% =================== zo werk je met dit boek ===================
\renewcommand\kopL{Voorwerk}\renewcommand\kopR{Zo werk je met dit boek}
\phantomsection\label{uitleg}\pdfbookmark[-1]{Zo werk je met dit boek}{uitleg}
{\sffamily\zwart\fontsize{17}{20}\selectfont Zo werk je met dit boek}\par\vskip 6pt
{\opmaakproza\small
Dit boek bevat echte examenopgaven <<VAKKLEIN>> <<NIVEAU>> uit de centrale examens van <<JAREN>>,
precies zoals ze op het examen stonden.

\textbf{Deel 1: oefenen per onderwerp.} <<NOPG>> opgaven uit <<NEXAMENS>> examens, ingedeeld naar de domeinen van de
syllabus. Elke opgave staat bij het onderwerp waar de meeste punten over gaan. Begin bij een onderwerp dat je
lastig vindt, of werk de hoofdstukken op volgorde door.

\textbf{Deel 2: complete examens.} De <<NEX>> nieuwste examens in hun geheel. Maak ze aan het eind van je voorbereiding,
in \'e\'en keer en met de klok erbij.

\textbf{Uitwerkingen.} Achterin staat bij elke vraag het beoordelingsmodel uit het offici\"ele correctievoorschrift:
de stappen (met \textbullet) en rechts hoeveel punten elke stap oplevert. Zo kijk je je werk na zoals een docent dat doet.
Een andere aanpak dan het model is ook goed als die <<VAKWOORD>> klopt; vaak staan er meerdere antwoordalternatieven.
Een \textcolor{kleur2}{\sffamily\scriptsize\bfseries AANVULLING} is een latere correctie van het CvTE.

\begin{tikzpicture}[x=1mm,y=1mm,font=\sffamily\scriptsize]
  \fill[papier] (0,0) rectangle (112,25);
  \node[anchor=north east,align=right,inner sep=0] at (14,21) {{\bfseries\fontsize{10.5}{11}\selectfont\color{kleur1}7}\\[-.5pt]{\fontsize{7}{8}\selectfont\color{grijs}4p}};
  \draw[grijs!60] (17,21.5) -- (17,13);
  \node[anchor=north west,text width=88mm,align=left,inner sep=0] at (20,21.5) {Het nummer van de vraag en het aantal punten (4p). In deel 1 begint de nummering bij elk onderwerp opnieuw, in deel 2 zijn het de nummers uit het examen. In de uitwerkingen staat hetzelfde nummer.};
  \node[anchor=north west,inner sep=0] at (3,9) {\colorbox{kleur1}{\color{white}\bfseries\fontsize{7.5}{9}\selectfont\strut 2.5}};
  \node[anchor=north west,text width=88mm,align=left,inner sep=0] at (20,9.2) {Opgave 5 van hoofdstuk 2. Rechts van de titel: uit welk examen de opgave komt.};
\end{tikzpicture}

\textbf{Tabs.} De gekleurde tabs aan de rand van de bladzijden laten zien in welk hoofdstuk je bent; de uitwerkingen
van een hoofdstuk hebben dezelfde kleur en dezelfde plek.

\textbf{Afwijkingen.} Opgaven en figuren zijn verkleind tot A5; de tekst is ongeveer 10 punt.
Waar in een vraag een uitwerkbijlage wordt genoemd, staat die direct onder de vraag.\par}
\clearpage

% =================== het examen en je cijfer ===================
\renewcommand\kopR{Het examen en je cijfer}
\phantomsection\label{cijfer}\pdfbookmark[-1]{Het examen en je cijfer}{cijfer}
{\sffamily\zwart\fontsize{17}{20}\selectfont Het examen en je cijfer}\par\vskip 6pt
{\opmaakproza\small
Het centraal examen <<VAKKLEIN>> <<NIVEAU>> duurt <<DUUR>>. Het bestaat uit ongeveer <<GEMVR>> vragen, verdeeld over
enkele opgaven, en je kunt er ongeveer <<GEMPT>> punten mee halen. Bij elke vraag staat hoeveel punten hij waard is.
Hulpmiddelen: <<HULP>>

\textbf{Van score naar cijfer.} Na het examen bepaalt het CvTE de normering (de N-term) en publiceert die op examenblad.nl. Je cijfer is
\[ \text{cijfer} = 9{,}0 \times \frac{\text{jouw score}}{\text{maximale score}} + N \]
met $N$ tussen 0,0 en 2,0; <<NTERM>>. Met $N=1$ heb je dus
ongeveer de helft van de punten nodig voor een 5,5. Bij een hoge of lage $N$ gelden nog extra regels, zodat
een score van 0 altijd een 1,0 en de maximale score altijd een 10 oplevert.

\textbf{Slagen.} Je eindcijfer voor het vak is het gemiddelde van je schoolexamencijfer en je cijfer voor het centraal examen.
Om te slagen moet onder meer het gemiddelde van al je centraal-examencijfers minstens een 5,5 zijn.

\textbf{Je tijd verdelen.} <<DUUR>> voor ongeveer <<GEMPT>> punten is <<MINPUNT>> minuten per punt. Een vraag van 4 punten mag dus
zo'n <<MIN4>> minuten kosten. Kom je er niet uit, sla de vraag dan over en kom er aan het eind op terug.

\textbf{Nakijken.} Hoe correctoren de punten geven, staat bij de uitwerkingen (blz.~\pageref{nakijken}).\par}
\clearpage

% =================== tips ===================
\renewcommand\kopR{Tips voor het examen}
\phantomsection\label{tips}\pdfbookmark[-1]{Tips voor het examen}{tips}
{\sffamily\zwart\fontsize{17}{20}\selectfont Tips voor het examen}\par\vskip 6pt
\input{tips}
\clearpage

\ifformulekaart
% =================== formulekaart ===================
\renewcommand\kopR{Formulekaart}
\phantomsection\label{formulekaart}\pdfbookmark[-1]{Formulekaart}{formulekaart}
{\sffamily\zwart\fontsize{17}{20}\selectfont Formulekaart}\par\vskip 4pt
{\small Deze formules staan voorin elk examen. Je mag ze tijdens het examen gebruiken.\par}\vskip 8pt
\input{formulekaart}
\clearpage
\fi

% =================== deel 1 ===================
\Deel{1}{Oefenen per onderwerp}{<<NOPG>> opgaven uit de centrale examens, ingedeeld naar de domeinen van de syllabus. De uitwerkingen staan achterin.}{kleur1}
\input{deel1}

% =================== deel 2 ===================
\Deel{2}{Complete examens}{De <<NEX>> nieuwste centrale examens <<VAKKLEIN>> <<NIVEAU>>, zoals ze zijn afgenomen. Maak ze in \'e\'en keer, in <<DUUR>>.}{kleurex}
\input{deel2}

% =================== uitwerkingen ===================
\Deel{3}{Uitwerkingen}{Bij elke vraag het beoordelingsmodel uit het offici\"ele correctievoorschrift van het CvTE, met de aanvullingen die later zijn gepubliceerd.}{kleur6}
\renewcommand\Tab{0}\renewcommand\kopL{Uitwerkingen}\renewcommand\kopR{Zo wordt er nagekeken}
\phantomsection\label{nakijken}\pdfbookmark[0]{Zo wordt er nagekeken}{nakijken}
{\sffamily\zwart\fontsize{15}{18}\selectfont Zo wordt er nagekeken}\par\vskip 5pt
{\opmaakproza\small Correctoren volgen het beoordelingsmodel: elke stap met een \textbullet\ levert de punten op die
erachter staan. Daarnaast gelden algemene regels: een fout wordt in een vraag maar \'e\'en keer aangerekend, en als je
meer antwoorden geeft dan gevraagd, telt alleen het eerste. Voor <<VAKKLEIN>> gelden bovendien deze regels uit het
correctievoorschrift van het nieuwste examen:\par}
\vskip 6pt
\input{regels}
\clearpage
\input{antw1}
\clearpage
\input{antw2}

% =================== nawerk ===================
\clearpage
\colorlet{hs}{kleur6}\renewcommand\Tab{0}\renewcommand\kopL{Overzicht}\renewcommand\kopR{Overzicht per examen}
\phantomsection\label{register}\pdfbookmark[-1]{Overzicht per examen}{register}
{\sffamily\zwart\fontsize{17}{20}\selectfont Overzicht per examen}\par\vskip 4pt
{\small Waar elke opgave in dit boek staat, per examen. Zo kun je ook een heel examen uit deel 1 maken.\par}
\input{register}
\clearpage
\renewcommand\kopL{Bronnen}\renewcommand\kopR{Bronvermelding}
\phantomsection\label{bronnen}\pdfbookmark[-1]{Bronvermelding}{bronnen}
{\sffamily\zwart\fontsize{17}{20}\selectfont Bronvermelding}\par\vskip 4pt
{\small Alle examenopgaven en uitwerkingen: College voor Toetsen en Examens (CvTE), examenblad.nl, \textcopyright\ Staat der Nederlanden.
Hieronder per examen de bronnen van de figuren en foto's, zoals vermeld in het correctievoorschrift, en de originele documenten.\par}
% rechts rafelig: bronnen bevatten lange codes en adressen die niet netjes uitvullen
{\raggedright\input{bronnen}\par}

% =================== succes ===================
% de laatste pagina voor de achterkant: de achterkant is een even pagina, dus deze een oneven
\clearpage
\ifodd\value{page}\else\thispagestyle{leeg}\mbox{}\clearpage\fi
\thispagestyle{leeg}\noindent
\phantomsection\pdfbookmark[-1]{Succes!}{succes}
\begin{tikzpicture}[remember picture,overlay]
  % confetti bovenin, vast patroon zodat elk boek hetzelfde is
  \foreach \x/\y/\r/\c in {22/22/20/kleur1,34/15/-30/kleur2,47/27/45/kleur3,61/12/10/kleur4,76/24/-15/kleur1,
      90/14/35/kleur2,104/26/-40/kleur3,118/16/25/kleur4,129/29/-10/kleur2,40/31/60/kleur4,84/32/-55/kleur3,112/33/15/kleur1}
    {\node[fill=\c,inner sep=0,minimum width=2.6mm,minimum height=4.4mm,rotate=\r] at ($(current page.north west)+(\x mm,-\y mm)$) {};}
\end{tikzpicture}\par
\vspace*{30mm}
{\sffamily\zwart\fontsize{26}{30}\selectfont Succes met je\\examen!\par}\vskip 6mm
{\opmaakproza\fontsize{11}{16.5}\selectfont
Je hebt geoefend met echte examenvragen <<VAKKLEIN>> uit <<JAREN>>. Wat er straks ook op je examen staat:
dit soort vragen heb je eerder gezien.\par}
\vskip 7mm
{\opmaakproza\small\setlength\parskip{1.6mm}
{\sffamily\bfseries\color{kleur1}De avond ervoor}\par
Leg je spullen klaar: pennen, potlood, gum en je legitimatie. Bij dit examen mag je ook gebruiken: <<HULPKLEIN>>
Leer niets nieuws meer, blader hooguit door de vragen die je fout had. En ga op tijd slapen.\par
\vskip 2mm
{\sffamily\bfseries\color{kleur1}Tijdens het examen}\par
Lees elke vraag rustig, en begin bij wat je kunt. Houd ongeveer <<MINPUNT>> minuten per punt aan.
Laat niets leeg: ook een begin van een antwoord kan punten opleveren.\par
\vskip 2mm
{\sffamily\bfseries\color{kleur1}Daarna}\par
Klaar is klaar. Kijk vooruit naar het volgende examen, niet terug naar dit.\par}
\vfill
{\sffamily\zwart\large\color{kleur2}Op naar je diploma!\par}
\vskip 2mm
{\sffamily\small\color{grijs}<<SAMENSTELLER>>\par}

% =================== achterkant ===================
\clearpage
\ifodd\value{page}\thispagestyle{leeg}\mbox{}\clearpage\fi
\thispagestyle{leeg}
\pdfbookmark[-1]{Achterkant}{achterkant}
\begin{tikzpicture}[remember picture,overlay]
  \definecolor{omslag}{HTML}{12313B}
  \fill[omslag] (current page.south west) rectangle (current page.north east);
  \begin{scope}[shift={(current page.south west)}]
    \begin{scope}[white,opacity=.07,line width=.4pt]
      \foreach \i in {0,...,42} {\draw (0,\i*5mm) -- (148mm,\i*5mm);}
      \foreach \i in {0,...,29} {\draw (\i*5mm,0) -- (\i*5mm,210mm);}
    \end{scope}
  \end{scope}
  \node[anchor=north west,text=white,align=left,text width=112mm] at ($(current page.north west)+(18mm,-30mm)$)
    {{\sffamily\zwart\fontsize{17}{21}\selectfont Oefenen met echte examens}\\[5mm]
     {\fontsize{10.5}{15}\selectfont Dit boek bevat <<NVR>> vragen uit de centrale examens <<VAKKLEIN>> <<NIVEAU>> van <<JAREN>>,
      ingedeeld naar onderwerp, plus <<NEX>> complete examens. Bij elke vraag staat het offici\"ele beoordelingsmodel van
      het College voor Toetsen en Examens, zodat je je werk precies zo kunt nakijken als op het examen.\\[4mm]
      De examens zijn openbaar en van iedereen. Daarom is dit boek gratis.}};
  \node[anchor=south west,text=white,font=\sffamily\scriptsize,align=left] at ($(current page.south west)+(18mm,14mm)$) {<<SAMENSTELLER>>\\ <<WEBADRES>>};
\end{tikzpicture}
\end{document}
